\chapter{Introduction} \label{chap:introduction}


The ability of cells to move is central to life.
During embryonic development, cells migrate over long distances to build tissues and organs~\cite{scarpa2016collective, weijer2009collective}; in the adult, they crawl to close wounds~\cite{poujade2007collective} and patrol the body as part of the immune response~\cite{trepat2012cell}; and in disease, the same migratory ability allows cancer cells to escape a primary tumor and seed metastases~\cite{haeger2014cell, cheung2016collective}.
Many of these processes are not carried out by isolated cells but by groups that move, grow, and remodel together~\cite{cheung2025collective, friedl2009collective, rorth2009collective}.
A recurring theme is that the behavior of such a group is more than the sum of its cells: collective motion, mechanical properties, and the ability to sense the environment all emerge from interactions between neighboring cells~\cite{alert2020physical, hakim2017collective}.

This dissertation studies two faces of this collective behavior.
The first is mechanical: how a group of cells holds together and, just as importantly, how it comes apart.
Tissues must be cohesive enough to maintain their shape, yet many biological processes depend on cells detaching in a controlled way~\cite{friedl2009collective, cheung2025collective}.
A primary tumor extends a strand of cells into the surrounding tissue and then sheds single cells and small clusters that can seed metastases~\cite{cheung2016collective, haeger2014cell}.
Germline cysts in the mouse reach different characteristic sizes during embryonic development through a balance between cell division, which grows the cluster, and fracture of intercellular bridges, which breaks it into smaller groups~\cite{levy2024tug}.
The metazoan \textit{Trichoplax adhaerens}, one of the simplest known animals, reproduces by a motility-induced fracture~\cite{prakash2021motility}.
In each case the outcome is set by mechanics: cells adhere to one another, crawl, generate forces, and eventually detach.
This raises a pair of physical questions that run through the first half of this dissertation: what controls when a group of cells breaks apart, and what sets the size of the pieces that result?

The second face is informational: to move sensibly, a cell must read its surroundings.
Cells follow chemical gradients~\cite{camley2016emergent, ipina2022collective}, mechanical cues~\cite{pallares2023stiffness}, and electromagnetic signals~\cite{sun2023electric, nwogbaga2023physical}, and they do so with surprising accuracy despite the noise inherent in counting a finite number of molecules~\cite{berg1977physics}.
A particularly basic version of this problem is sensing other cells.
When two crawling cells collide, each often repolarizes and the two move apart, a behavior known as contact inhibition of locomotion (CIL)~\cite{abercrombie1970contact, mayor2010keeping}, which can be triggered by a contact only micron-scale in size~\cite{singh2021rules}.
This motivates the question taken up in the second half of the dissertation: how accurately can a cell sense that it has touched a neighbor?

This dissertation develops theoretical and computational models for these two problems: the mechanics of how cell groups detach and fracture, and the physical limits on how a cell senses its neighbors.
My aim throughout is to use simple models to expose the physical mechanisms behind these observed behaviors.


\section{Insights from experiments}
\label{sec:ch1:experiments}

Collective cell migration appears in many biological processes, from embryonic development to wound healing to cancer~\cite{friedl2009collective, cheung2025collective}.
The border cells of the \textit{Drosophila} egg chamber move as a small cluster of six to eight cells through the surrounding nurse cells~\cite{montell2012group}; the vertebrate neural crest migrates as multicellular streams that disperse across the embryo to form cartilage, peripheral nerves, and pigment cells~\cite{scarpa2016collective}; and during wound healing a single epithelial sheet advances over the wound bed to close it~\cite{poujade2007collective}.
In cancer, invasion of the surrounding tissue is often collective, with strands and buds of tumor cells extending from the primary tumor mass~\cite{cheung2025collective, friedl2009collective}.
In each case the cells stay attached to their neighbors through cadherin junctions and coordinate their motion, so that the group migrates as one mechanical unit~\cite{friedl2009collective}.

A group of cells can also display collective properties not seen in single cells~\cite{cheung2025collective}.
A confluent monolayer can jam into a solid-like state, in which each cell is caged by its neighbors, or unjam into a fluid-like state, in which cells rearrange and flow past one another, with the transition controlled by cell shape, motility, and adhesion~\cite{park2015unjamming, bi2016motility, rustarazo2026adhesion}, and this has been proposed as a route by which tumor cells escape a primary tumor and invade~\cite{cai2025unjamming}.
A migrating group can also be polarized supracellularly, as a whole rather than cell by cell: the cells at the front lead with forward protrusions while the rear is contractile, so the group has a single front-to-back axis, much like one large cell~\cite{friedl2009collective, cheung2025collective}.
For instance, in the cranial neural crest of \textit{Xenopus}, the rear cells build a shared actomyosin cable that runs along the back of the cohort and helps propel it forward~\cite{cheung2025collective}.

A question that recurs across these systems is how a group stays cohesive and yet still lets cells leave in a controlled way~\cite{friedl2009collective, cheung2025collective}.
One route is the epithelial-mesenchymal transition (EMT), in which epithelial cells loosen their junctions and become individually migratory, a program that is reactivated during cancer invasion~\cite{youssef2024epithelial, tripathi2020physics}.
But cells need not complete that transition to leave: in many epithelial cancers the disseminating cells retain E-cadherin and their cell-cell junctions, so cohesive clusters break away and travel together~\cite{cheung2025collective}.
What sets when a group fragments, and how large the resulting pieces are, is then a question of mechanics, of how cells adhere, crawl, and pull on one another.

\textit{In vivo}, collective migration is often confined by adjacent tissues, by the dense fibers of the extracellular matrix (ECM), and by the vasculature, the walls of blood and lymphatic vessels~\cite{tarle2017modeling, balzer2012physical, lin2025gradients}.
During cancer invasion, the primary tumor extends a strand of cells into the surrounding tissue, and intravital microscopy shows that tumor cells preferentially move along channel-like tracks created by these anatomical structures, with widths from a few microns up to wider than a single cell~\cite{paul2017cancer, weigelin2012intravital, wolf2009collagen}.
Whether cells leave the strand one at a time or as clusters matters for the disease: clusters of circulating tumor cells are rare, but they seed metastases far more efficiently than single cells~\cite{cheung2016collective, green2022pillarx, dirusso2024beyond}.
Moreover, it has been reported that circulating tumor cells (CTCs) may even cluster with white blood cells or platelets to escape immune surveillance~\cite{saini2019circulating}.
To re-create this confinement \textit{in vitro}, Law \textit{et al.} induced human A431 carcinoma cells by a serum gradient to invade microchannels of different widths, from $6$ to $50\,\mu\mathrm{m}$, and tracked single cells and clusters detaching from the invading front~\cite{law2023cytokinesis}.
Most detachment events involve a single cell, and larger clusters (up to roughly twenty cells in the widest channels) detach more often in wider channels.
Surprisingly, the time a strand takes to rupture does not depend on the channel width.
These observations raise the question of what mechanical features of the strand and the channel set the rupture statistics, and motivate the model developed in Chapter~\ref{chap:dissociation}.

Recent experiments on germline cysts in the mouse have shown that two competing processes can set the size of a small cell cluster: cell division builds the cluster, while random cell motility fractures the intercellular bridges that connect its cells~\cite{levy2024tug}.
A related mechanism drives reproduction in the metazoan \textit{Trichoplax adhaerens}, one of the simplest known animals.
\textit{Trichoplax} has no nervous system, no basement membrane, and no visible extracellular matrix; its body is a flat sheet of two epithelial layers only about $25\,\mu\mathrm{m}$ thick but several millimeters wide~\cite{prakash2021motility, ueda1999dynamic, smith2019coherent}.
This minimal anatomy makes it simple enough to model and yet still a complete animal, and it reproduces by motility-induced fracture as it crawls~\cite{prakash2021motility}.
A similar fracture-limited size control occurs in multicellular snowflake yeast, though there it is driven by crowding stress rather than by motility~\cite{jacobeen2018cellular}.
Tissue and organ sizes are more commonly thought to be set by regulation of cell proliferation~\cite{lecuit2007orchestrating}, but in these systems they instead emerge from a balance of growth and fracture.
Moreover, both \textit{Trichoplax} and germline cysts show significant variability in sizes despite the presence of a characteristic size scale~\cite{prakash2021motility, levy2024tug, ikami2023branched, lei2013mouse}.
This combination of a well-defined typical size with substantial individual-to-individual variability raises the question of how reliably motility-driven fracture can regulate size at all.
This question is taken up in Chapter~\ref{chap:size-control}.

Contact between cells does more than hold a group together: through contact inhibition of locomotion (CIL), it also helps steer and organize collective migration.
When two crawling cells collide, each retracts its protrusions at the point of contact, repolarizes, and the two move apart, a behavior first seen in fibroblasts by Abercrombie more than half a century ago~\cite{abercrombie1970contact, abercrombie1979contact} and since found in many healthy and cancerous cell types~\cite{mayor2010keeping, zimmermann2016contact}.
Because each cell turns away from the neighbors it touches, contacts in the interior of a group suppress protrusions there and bias them toward the free edge, so that the group spreads outward or streams in a common direction; in the vertebrate neural crest, this contact-driven repolarization controls the directional migration of the cell streams~\cite{scarpa2016collective, mayor2010keeping}.
CIL also sets how cells space themselves out, and \textit{Drosophila} macrophages rely on it to disperse evenly through the embryo~\cite{davis2012emergence, stramer2010clasp}.
At the molecular level the response can be triggered by the Eph-ephrin pathway, in which Eph receptors on one cell bind ephrin ligands anchored to the surface of the cell it touches~\cite{larranaga2022ephrin, lin2015interplay}, and it is precise enough to fire even when two cells crawling on nanofibers touch over a patch only a micron wide~\cite{singh2021rules}.
Detecting a neighbor this way can therefore be cast as a sensing problem, namely how reliably a cell can tell that a neighbor is present, which we study in Chapter~\ref{chap:cil}.

\section{Physical models of collective cell migration}
\label{sec:ch1:models}

\subsection{Agent-based models of tissues}

\begin{figure}[tbp]
\centering\vspace{0.3cm}
\includegraphics[width=\textwidth]{chap1/models.pdf}
\vspace{0.3cm}
\caption{(a) Examples of common models of collective cell migration. Reprinted from Fig.~2 of Ref.~\cite{camley2017physical}, with permission from IOP Publishing (Order License ID: 1739114-1). (b) A schematic sketch of proteins associated with front (in shades of red) and rear (green) of a polarized cell. Reprinted from Fig.~1D of Ref.~\cite{rappel2017mechanisms}, with permission from Elsevier (License Number: 6285460585190). (c) The structure of the Eph receptors and their ephrin ligands. Reprinted from Fig.~1 of Ref.~\cite{coulthard2012eph}, under the terms of the CC BY-NC-ND 4.0 license.}
\label{fig:ch1:models}
\end{figure}

The molecular biology of cell motility is rich: integrins coupling the actin cytoskeleton to the extracellular matrix, E-cadherins coupling cells to each other, Rho-family GTPases setting the polarity, and branched actin networks pushing the cell forward~\cite{parsons2010cell, ridley2003cell, demali2003coupling, dufour2013alpha, rappel2017mechanisms}.
A faithful model of this molecular detail would itself be complex.
For the tissue-scale questions of this dissertation, however, such a faithful description is usually not necessary.
What matters is whether a few simple physical ingredients, mechanical contact, an active velocity, a few rate constants, are enough to reproduce a given experimental observation.
Physical models of motile cells therefore vary widely in their level of detail, from cells as single self-propelled points to fields that resolve each cell's shape [Fig.~\ref{fig:ch1:models}(a)], and each chapter of this dissertation uses a model from this hierarchy chosen to match its question.
A more complete survey of the modeling toolbox is given by the review~\cite{camley2017physical}.

The simplest models treat each cell as a self-propelled particle (SPP)~\cite{fily2012athermal, levine2000self, copenhagen2018frustration, sanoria2026chemotaxis}.
In the active Brownian particle (ABP) model the velocity has fixed magnitude and a direction that rotates diffusively, while in the active Ornstein-Uhlenbeck particle (AOUP) model the velocity itself is an Ornstein-Uhlenbeck (OU) process with fluctuating magnitude and direction~\cite{dunn1987unified, fily2012athermal, fodor2016far}.
These models are cheap enough to simulate many particles and are often analytically tractable, which makes them well suited to statistical questions about clustering, fracture rates, and steady-state distributions~\cite{fodor2016far, henkes2020dense, wang2026controlling}.
Chapter~\ref{chap:size-control} uses the active Ornstein-Uhlenbeck model to study tissue size control by motility-driven fracture.
What such models give up is cell shape. Moreover, dense packings of identical particles tend to crystallize into ordered lattices~\cite{prakash2021motility}, which is a problem for modeling real tissues, where the cells are disordered.

Adding cell shape gives vertex and Voronoi models for confluent tissue, in which cells tile space without gaps~\cite{bi2015density, bi2016motility, henkes2020dense, park2015unjamming}.
In a vertex model, each cell is a polygon defined by its vertices, and the energy penalizes departures from a preferred area and perimeter~\cite{fletcher2014vertex}.
Because the vertices outnumber the cells, the number of degrees of freedom is larger than in a particle model, and the model predicts a rigidity transition tuned by the cell shape index, with floppy modes near the transition~\cite{bi2015density, park2015unjamming, lawson2021jamming}.
The active Voronoi model is more economical: each cell is the Voronoi region of its center, so the only degrees of freedom are the $N$ cell-center positions, the same count as a particle model.
This geometric constraint matters.
At zero temperature the static Voronoi model lacks the shape-index rigidity transition seen in vertex models because each cell's shape is fixed by the Voronoi tessellation of the cell centers, leaving no independent shape degrees of freedom to tune~\cite{sussman2018no}; the active Voronoi model with self-propulsion does, however, exhibit a motility-driven glass transition~\cite{bi2016motility}.
However, because cells in either model tile space, neither allows cells to pull apart.

Phase-field models go further, representing each cell as a smooth field on a grid that is one inside the cell and zero outside~\cite{camley2014polarity, palmieri2015multiple, loewe2020solid, perez2024deposited, zadeh2024inferring, chiang2024multiphase}.
The dynamics of these fields are motivated by classical phase-separation equations.
The Cahn-Hilliard equation describes a strictly conserved order parameter and is more commonly applied to systems such as liquid droplets with fixed total volume~\cite{cahn1958free, negro2023yield}.
The Allen-Cahn equation describes an unconserved order parameter and is the more natural starting point for cells, since the cell-interior indicator is not itself a conserved quantity; cell area is instead held near a preferred value by a soft energetic penalty, in the same spirit as the area term in vertex and Voronoi models~\cite{loewe2020solid, negro2023yield}.
Because the cell boundary is free to take any shape, phase-field models resolve deformation, allow cells to separate, and reproduce the unjamming transition of confluent tissues~\cite{loewe2020solid, chiang2024multiphase}.
The cost is computational: representing every cell as a field on a grid limits phase-field simulations to relatively small numbers of cells.
Chapter~\ref{chap:dissociation} uses a phase-field model for cancer cell dissociation, where resolving cell shape and detachment matters more than reaching large system sizes.

For larger nonconfluent tissues, such as the metazoan \textit{Trichoplax adhaerens} (an animal composed of up to $\sim 10^5$ cells), neither family of models is ideal: phase-field simulations do not scale to that many cells, and the standard vertex and Voronoi models tile space without gaps and so cannot represent a nonconfluent tissue.
Some intermediate approaches have been proposed to bridge the two regimes.
The finite Voronoi model truncates each cell's Voronoi region at a maximum radius, so that cells keep a polygonal shape where they touch but can also pull apart and leave gaps between cells~\cite{graner1993can, teomy2018confluent, huang2023bridging}.
Deformable polygon models instead represent each cell by an explicit polygon with many vertices, without imposing a Voronoi geometry~\cite{boromand2018jamming, lv2024active}.
Subcellular element models go further by representing each cell as a collection of interacting subunits~\cite{sandersius2008modeling}.
Chapter~\ref{chap:voronoi} studies the finite Voronoi model in detail.


\subsection{Modelling interactions and polarity dynamics}
\label{sec:ch1:interactions}

A crawling cell anchors itself to its environment through transmembrane integrin receptors, which couple the extracellular matrix (ECM) to the actin cytoskeleton at focal adhesions and generate the traction force needed to crawl~\cite{parsons2010cell, ananthakrishnan2007forces}.
Two cells in contact form a different kind of adhesion: cadherins on the membrane of one cell bind cadherins on the other, and the cytoplasmic tails of these cadherins are linked through catenins to the actin cortex, mechanically coupling the two cells~\cite{noordstra2023cadherin, dufour2013alpha}.
A third class of contacts uses receptor tyrosine kinases such as Eph receptors that bind ephrin ligands on the neighbor, triggering downstream signaling rather than a long-lived mechanical attachment~\cite{larranaga2022ephrin, lin2015interplay}.

Physical models capture these interactions through effective potentials and energies that summarize the underlying molecular and mechanical detail.
For cell-cell coupling, particle models use a short-range pairwise potential with soft repulsion and an attractive tail~\cite{sanoria2026chemotaxis, copenhagen2018frustration}, or explicit breakable spring junctions as in the chain of Chapter~\ref{chap:size-control}; vertex and Voronoi models encode adhesion in a line tension on cell-cell interfaces, where a lower tension corresponds to stronger adhesion and a longer preferred contact~\cite{bi2015density, bi2016motility, huang2023bridging}; and phase-field models use an interaction energy between the fields of neighboring cells, typically a surface-overlap term that lowers the energy when two cell boundaries coincide~\cite{zadeh2024inferring, camley2014polarity}.
Cell-environment coupling enters analogously, as a soft-wall potential or an attraction to a substrate field.
For numerical stability at strong adhesion, Chapter~\ref{chap:dissociation} adapts the phase-field framework by introducing an advection term that handles both the cell-cell and cell-wall adhesion.

The direction a cell wants to go is set by its polarity, an asymmetry between leading and trailing edges that is maintained by the Rho family of small GTPases~\cite{rappel2017mechanisms}.
As sketched in Fig.~\ref{fig:ch1:models}(b), Rac1 and Cdc42 are active at the leading edge, where they activate the Arp2/3 complex to nucleate the branched actin networks that push out the lamellipodium~\cite{demali2003coupling, ananthakrishnan2007forces}.
RhoA is active at the rear, where it promotes the actomyosin contractility that pulls the cell body~\cite{rappel2017mechanisms, ananthakrishnan2007forces}.
External cues, including chemical gradients, electromagnetic fields, the mechanical stiffness of the substrate, and physical contact with a neighbor, reshape this polarity by changing the spatial distribution of these GTPases~\cite{rappel2017mechanisms, demali2003coupling, nwogbaga2023physical, klein2012eph}.

Physical models usually reduce this machinery to a phenomenological dynamics for the cell's polarity vector, the direction in which it tends to move.
In the simplest case the polarity diffuses freely, giving the persistent random walks of active Brownian and active Ornstein-Uhlenbeck particles introduced above.
The Vicsek model is the canonical example of velocity alignment: each particle rotates its velocity toward the average velocity of its neighbors, originally for flocks of self-propelled particles and later for cells~\cite{vicsek1995novel}.
A common way to model chemotaxis is to bias the polarity along an external chemical gradient, while in models of CIL, the polarity is most often rotated away from contacting neighbors, taken as the vector sum of unit vectors pointing away from each contact~\cite{camley2016emergent, ipina2022collective, sanoria2026chemotaxis, copenhagen2018frustration}.
The phase-field model of Chapter~\ref{chap:dissociation} includes guided cells whose polarity vector points along the serum gradient imposed in the experiments.


\section{Physical limits on cell sensing}
\label{sec:ch1:sensing}

To migrate purposefully, adopt the correct developmental fate, and respond to neighboring cells, cells must sense and interpret their surroundings despite noisy molecular signals~\cite{sengupta2021principles}.
These functions can all be cast as problems of information sensing and processing: chemotaxing cells follow shallow chemical gradients~\cite{hu2010physical}, cells in a developing tissue interpret morphogen concentrations to decide their fate~\cite{rogers2011morphogen}, and immune cells discriminate self from non-self at a cell-cell contact~\cite{feinerman2008quantitative}.
A central question in cell biophysics is how accurately a cell can perform such tasks given that the underlying signals, the arrival of individual molecules at individual receptors, are inherently stochastic~\cite{berg1977physics}.
Answering it sets a benchmark: If an observed behavior approaches the physical limit, the cell is using its information nearly optimally, whereas a large gap points to additional noise or unused information.

A cell measures a chemical signal by counting molecules that bind to receptors on its membrane.
Berg and Purcell showed that this counting is itself a fundamental source of noise: because molecules arrive at random, a receptor monitored for a finite time returns only an estimate of the concentration, with an error that shrinks slowly as the observation time and the number of receptors grow~\cite{berg1977physics}.
This argument sets a physical limit that no downstream processing can beat, and it has become the starting point for a large body of theory on the physics of cellular information processing~\cite{bialek2005physical, ten2016fundamental, tkavcik2025information}.

The Berg-Purcell argument has been extended to the harder tasks that real cells perform.
Cells do not only measure a uniform concentration; they need to estimate signals that change in space and time.
They read spatial gradients to decide which way to crawl~\cite{hu2011how, hu2010physical}, they estimate the position of a spatially localized signal~\cite{wasnik2019accuracy}, they sense gradients collectively as a group~\cite{ipina2022collective, camley2016emergent}, and they track signals that change in time~\cite{mora2010limits}.
Even for the single task of finding a gradient, a cell faces a choice of strategy: it can compare concentrations across its body at one instant (spatial sensing) or compare them at different points along its path over time (temporal sensing), and which is more accurate depends on the cell's size, speed, and persistence~\cite{rode2024information}.
The cell's shape matters too: because gradient information is limited by how widely the receptors are dispersed, an elongated or branched shape can sense a shallow gradient more accurately than a round one~\cite{mou2025optimal}.
A useful way to treat all of these problems is as statistical estimation: the cell is assumed to extract as much information as it can from its receptors using maximum-likelihood estimation (MLE), and the best accuracy any such estimate can reach is bounded by the Cram\'er-Rao bound~\cite{endres2009maximum}.
This framework gives a clean way to ask what limits a given sensing task, separate from the details of any particular signaling network.

Several broad questions have guided this field.
A major line of research concerns the fundamental limits and asks how cells can beat the naive Berg-Purcell limit: Time-averaging, spatial averaging over many receptors, and comparisons across the cell body can all reduce sensing error, although each strategy has its own trade-offs and ultimate limits.~\cite{ten2016fundamental}.
A related line of work asks how groups of cells sense better than individuals, sharing information across cell-cell contacts to detect gradients that are too shallow for a single cell to resolve~\cite{ipina2022collective, camley2016emergent}.
Some other questions concern adaptation and dynamic range, namely how a cell can stay sensitive to relative changes in signal across many orders of magnitude of background concentration~\cite{endres2006precise, srinivasan2026perfect}.

The second half of this dissertation focuses on a different but closely related problem: specificity.
Receptors are not perfectly specific.
A receptor tuned to one ligand will also bind other molecules that resemble it, and these spurious binding events add noise that is distinct from counting noise~\cite{mora2015physical, singh2017simple}.
When the correct ligand is rare and the interfering ligands are common, telling them apart becomes its own limit on sensing accuracy.
This problem arises whenever a cell must read a specific molecular cue against a noisy chemical background, from immune recognition to the detection of a neighboring cell.

CIL can be viewed in exactly these terms.
Research has shown that CIL is regulated by the Eph-ephrin pathway on the cell membrane [Fig.~\ref{fig:ch1:models}(c)]: Eph receptors on one cell bind ephrin ligands on the other and initiate the contact response, modulating the Rho GTPases of Sec.~\ref{sec:ch1:interactions} to increase actomyosin contractility and inhibit lamellipodium formation at the contact, so that the cell repolarizes and moves away~\cite{klein2012eph, noren2004eph, pasquale2008eph, arvanitis2008eph}.
For this response to fire reliably, a cell must first detect the presence and position of the cell-cell contact, a spatially localized signal made of ephrin ligands bound to Eph receptors~\cite{larranaga2022ephrin, lin2015interplay}.
The detection is made harder by interfering ligands that bind the same receptors~\cite{wykosky2008soluble, gong2016exosomes, minami2011ephrina}, and harder still when the contact is only micron-scale~\cite{singh2021rules}.
Chapter~\ref{chap:cil} treats CIL as a sensing problem of this kind, and asks how accurately a cell can locate a contact and how reliably it can decide that a neighbor is present at all.


\section{Overview of dissertation}
\label{sec:ch1:overview}

The remaining chapters are organized around the two main themes of my Ph.D. research.
The first three develop models of tissue rupture; the last develops a model of cell sensing.
Each is adapted from a published paper or a preprint, cited at the start of the chapter.

Chapter~\ref{chap:dissociation} studies how cancer cells dissociate from a collectively invading strand under confinement.
With a phase-field model of cells crawling through microchannels, I recapitulate the dissociation statistics measured by Law \textit{et al.} and find that single cells dissociate most often, that fast leader cells explain this predominance, and that stronger cell-cell adhesion produces fewer but larger ruptures.
Unjamming the tissue turns out to be necessary but not sufficient to create ruptures.

Building on this picture of fracture in a confined strand, Chapter~\ref{chap:size-control} asks whether motility-driven fracture can also regulate the size of a growing tissue.
Modeling a cluster as active particles joined by breakable junctions, I derive the exact steady-state size distribution and show that it depends only on the ratio of the junction break rate to the cell division rate.
Size control improves if division is restricted to the cluster edge or if fracture is localized to its middle.

To follow these results into larger nonconfluent tissues, such as \textit{Trichoplax adhaerens}, a more scalable model than phase-field is needed, one that still allows the tissue to remain fluid.
Chapter~\ref{chap:voronoi} examines the finite Voronoi model as a candidate, and shows that its detachment forces diverge as cells lose contact, so that the simulated time of fracture depends on the numerical time step rather than on the physics.
I then propose a regularization, calibrated against a deformable polygon model that shares the same energy function but does not impose a Voronoi shape.

Chapter~\ref{chap:cil} turns from mechanics to sensing and treats CIL as a problem of cell sensing.
Modeling a cell that detects a neighbor by binding ephrin to its Eph receptors against nonspecific interfering ligands, I use maximum-likelihood estimation and the Cram\'er-Rao bound to find the best possible sensing accuracy.
I find that a cell can decide whether a neighbor is present with high reliability even when the concentration of interfering ligands is ten times that of the correct ligand, and that the accuracy of locating the contact depends only weakly on its width.

Chapter~\ref{chap:conclusion} provides a summary of the dissertation and outlines potential directions for future research.
